% Chapter 4 — Experimental setup.
\section{Experimental Setup}

\subsection{Datasets}

The QA subset of HaluEval \cite{haluEval} is the training source. It contains 10{,}000 questions, each with a knowledge passage, a correct answer, and a hallucinated answer, which gives 20{,}000 labeled rows. The rows are split 70/15/15 at the question-group level, producing 14{,}000 training, 3{,}000 validation, and 3{,}000 test rows. Transfer is measured on RAGTruth \cite{ragtruth}, a corpus of naturally generated responses with human word-level span annotations, and on FaithBench \cite{faithbench}, a summarization hallucination benchmark built from modern LLM outputs. The RAGTruth QA task contributes 5{,}934 rows in 989 groups, of which 5{,}034 rows (839 groups) form the recalibration set and 900 rows (150 groups) the held-out test, with no overlapping groups between the two. Table~\ref{tab:datasets} summarizes the corpora.

\begin{table}[t]
\centering
\small
\caption{Datasets used for training and evaluation.}
\label{tab:datasets}
\begin{tabular}{@{}lcccl@{}}
\toprule
\textbf{Corpus} & \textbf{Rows} & \textbf{Groups} & \textbf{Use} \\
\midrule
HaluEval QA      & 20{,}000 & 10{,}000 & train / validation / test \\
RAGTruth QA      & 5{,}934   & 989      & zero-shot + recalibration \\
RAGTruth full    & 17{,}790  & 2{,}965   & zero-shot \\
FaithBench       & 750       & 750      & zero-shot \\
\bottomrule
\end{tabular}
\end{table}

\subsection{Evaluation protocol}

Models are compared on precision, recall, F1, AUROC, PR-AUC, and MCC, and calibration with expected calibration error over 10 bins, Brier score, and negative log-likelihood. Hyperparameters are chosen by a randomized search of 30 iterations with grouped 5-fold cross-validation on the training split. Every experiment is repeated with seeds 42, 123, and 456 and reported as the mean over seeds. Significance is assessed with McNemar's test on paired test predictions, bootstrap 95\% confidence intervals from 1{,}000 resamples, and a Wilcoxon signed-rank test across seeds. The group structure is respected everywhere, so the two rows of one question never appear in both the training and the evaluation folds.

\subsection{Calibration protocol}

Calibrators are fit on the validation split only, and the comparison of raw, Platt-scaled, and isotonic probabilities happens on the test split. The target-domain protocol follows the same principle. The recalibration set of 5{,}034 RAGTruth QA rows is used only to fit the target calibrator, and all reported numbers come from the disjoint 900-row test set.

\subsection{Explainability protocol}

Global explanations come from a SHAP TreeExplainer \cite{shap} over the raw XGBoost model. Three measurements check whether the explanations are trustworthy. First, the mean absolute SHAP ranking is compared against permutation importance with a Kendall tau correlation. Second, the top-5 feature set is re-estimated from 1{,}000 bootstrap resamples and compared with the Jaccard index. Third, two perturbation studies are run on the test answers. Content edits replace entities, numbers, or dates, remove the supporting context, or insert an irrelevant sentence, and each edit should move the calibrated score if it changes the evidence. Structure edits shuffle the clause order, and these should leave the score nearly unchanged. For each operator the mean absolute score change, the fraction of large score changes, the top-1 SHAP flip rate, and the Spearman correlation between the original and perturbed scores are recorded. A neutralization study sets the top-$k$ SHAP features of each sample to its median and tracks the mean score change for $k = 1, 3, 5, 10$.

\subsection{Efficiency and judge protocol}

Latency is measured on 200 test samples on the reference laptop with an RTX 3060 6\,GB GPU, where NLI and embeddings run on the GPU and the remaining features on the CPU. The LLM-as-judge comparison uses GPT 5.6 Luna on the same 200 samples with a fixed prompt, and the cost is computed from the recorded token counts.
